\selectlanguage{english}\Japanese
\Language{ja}{日本語}{利用ガイド}
\firsttopic{起動・スリープ・ノートへの復帰}
Inkline は reMarkable 2 上でローカルのシェルを動かします。Type Folio または USB
キーボードを使い、端末上で作業したり、SSH や Goblin Mosh で別のコンピューターに
接続したりできます。文字と画像は電子ペーパー向けのグレースケールで表示されます。

\begin{notice}
\key{Ctrl+Alt+T} で本体から Inkline を起動します。\par
\key{電源ボタン}を押して離すとスリープし、もう一度押すと復帰します。
開いているシェルやエディターのバッファーは RAM に保持されます。
\end{notice}
スリープは Inkline を終了せず、設定をフラッシュに保存しません。復帰後はネットワークの
再接続が必要な場合があります。復帰時のボタン操作では再びスリープしないようにしています。
更新後は一度 Inkline を終了して開き直すと、新しい電源ボタンの処理が有効になります。

ノートに戻るには \key{Quit} をタップするか \key{Ctrl+Shift+Q} を押し、確認します。
\code{exit} は現在のターミナルだけを閉じます。最後のターミナルを閉じるとノートに戻ります。
下部バーには Escape、履歴の上下、Quit、Settings があります。
起動時の短い案内と小さな Goblin ロゴは \code{clear} で消せます。

ショートカットアプリが画面を使用している場合は \key{Ctrl+Alt+T} で戻ります。
緊急時は \key{Ctrl+Alt+Backspace を2秒間}押し続けます。
この復旧操作は全ターミナルを閉じるため、セッション内の未保存の作業は失われます。
本体を再起動すると、通常どおりノート画面が開きます。

\ProjectLinks

\topic{キーとターミナルのスロット}
Type Folio の Ctrl と Alt の間にある独立した \key{Opt} キーと、
\key{右 Alt/Opt} キーは別の役割を持ちます。

\begin{keytable}
\key{キー} & \key{動作} \\ \midrule
独立した Opt + 1--0 & F1--F10 \\
右 Alt/Opt + 1--9 & スロット1--9を選択 \\
右 Alt/Opt + Left / Right & 前／次のスロット \\
右 Alt/Opt + Tab & Escape \\
右 Alt/Opt + Up / Down & PageUp / PageDown \\
右 Alt/Opt + Backspace & 全ターミナルの終了を確認 \\
左右どちらかの Alt + Space & Unicode キーボード \\
右 Alt/Opt + C / V & 選択範囲をコピー／貼り付け \\
Ctrl+Shift+B & 下部バーの表示／非表示 \\
\end{keytable}

\key{US 配列の Type Folio} では、右 Alt/Opt+\key{0} で \code{+}、
右 Alt/Opt+\key{Backspace のすぐ左にあるマイナスキー}で \code{=} を入力します。
入力方式が Off のとき、Shift+6 はそのままキャレットを入力します。
\code{c^2} は \code{c^2} のままです。文字サイズ変更のキー操作はありません。

USB キーボードの通常のファンクションキーも使えます。独立した Opt の操作は、
USB キーボードに Meta キーがある場合、その修飾キーを使います。

独立したスロットは最大 \key{9個}です。空のスロットを選ぶとシェルが開きます。
表示の分母は上限ではなく、現在開いている数です。1と9だけが開いていれば、
9の表示は \key{Terminal 2 / 2 · slot 9} です。閉じても他のスロット番号は変わりません。

各スロットには専用のシェル、作業ディレクトリー、変換中の入力、画像、500行の履歴があります。
表示と入力方式の設定は全スロットで共通です。

\topic{設定と画面の調整}
下部の5番目のボタン \key{Settings} をタップします。バーが非表示なら
\key{Alt+Space}、続いて \key{F2}、または Unicode キーボードの Settings を使います。
塗りつぶされた項目が現在の設定で、破線の枠はキーボードのフォーカスです。
Tab で移動し、矢印と Enter で変更します。

\key{Caps Lock：} Control（既定）と Caps Lock を切り替えられます。
\key{文字サイズ：} 2本指でピンチするか、設定のマイナス／プラスを使います。
ゆっくり動かすと1ピクセルずつ、\key{6～48ピクセル}の範囲で調整できます。

\key{Text darkness：} 50\% が元の描画です。値を上げると文字の縁が濃くなります。
\key{Minimum contrast} は、アプリが似た色を選んだときに前景と背景の明るさを離します。
既定値は35\%です。

\begin{choices}
\key{更新モード} & \key{特徴} \\ \midrule
Fast & 出力の待ち時間が短い。アニメーション用波形で残像は多め \\
Balanced & UI 用波形。32 ms 単位で出力をまとめる \\
Crisp（既定） & きれいな階調。コンテンツ用波形で更新は遅め \\
Mono & 高速な白黒表示。画像の中間階調は省略 \\
Saver & 120 ms 単位でまとめ、連続出力の更新回数を減らす \\
\end{choices}

変更された文字の行だけを再描画しますが、パネル自体の更新時間は残ります。
Saver の節電効果は未測定です。更新時も、明示的に保存したモードは維持されます。

\begin{notice}
設定は使用中 \key{RAM} に保持し、Inkline の通常終了時にまとめて保存します。
指を離す、設定を閉じる、スリープする操作では書き込みません。変更がなければ書き込みはゼロです。
クラッシュ、強制終了、電源断では未保存の変更が失われる場合があります。
\end{notice}

\topic{Unicode と入力方式}
\key{Alt+Space} で Unicode キーボードを開き、文字をタップして入力します。
連続入力のため画面は開いたままです。\key{Done} または Escape で閉じます。

文字、結合記号、数字、句読点、数学／通貨、記号、空白、書式制御、私用領域に分類されています。
上下にドラッグしてスクロールします。矢印、PageUp/PageDown、Home/End、マウスホイールも使えます。
\key{Tab} で分類を変更し、Enter で選択文字を入力します。

\code{03bb} のような16進の Unicode 値を入力して Enter を押すと λ が入ります。
Backspace で値を編集します。結合記号の点線の円は見本で、入力されるのは記号だけです。
空白と書式制御文字にはラベルがあります。文字一覧は Qt の Unicode データに基づき、
制御文字、未割り当て、サロゲート、非文字を除きます。未搭載のフォントが必要な文字もあります。
私用領域の表示には対応するフォントが必要です。

\key{F2} または \key{Settings} から \key{Input method} を選びます。
\begin{choices}
Off — direct keyboard & 通常の US 配列による直接入力 \\
Romaji & ローマ字から日本語のひらがなへ \\
Pinyin & ピンインによる中国語入力 \\
Zhuyin & 基本的な注音配列による中国語入力 \\
Wubi 86 & 字形コードによる中国語入力 \\
US-International & よく使うラテン文字のアクセント \\
\end{choices}

候補バーで結果を選びます。通常の US 入力に戻すには、同じメニューで
\key{Off — direct keyboard} を選びます。全ターミナルに適用され、終了時に保存されます。

\topic{タッチ・ペン・クリップボード}
\key{2本指で上下：} 履歴をスクロールします。下にドラッグすると古い出力、上に動かすと
プロンプトへ戻ります。各ターミナルは、現在の画面とは別に物理行500行を RAM に保持します。

\key{2本指で左右：} 1回のジェスチャーでスロットを1つ切り替えます。
次の切り替え前に指を離してください。指を広げたり狭めたりすると文字サイズが変わります。

\key{1本指：} ドラッグして文字を選択します。右 Alt/Opt+C でコピー、右 Alt/Opt+V で
貼り付けます。全ターミナルが最大128 KiB の RAM クリップボードを共有します。

\key{ペン：} マウス報告を有効にしたアプリには、ペンの押下、移動、離す操作を
左マウスのイベントとして送ります。SGR など、アプリが選んだプロトコルを使います。
それ以外では文字選択になります。\key{Shift} を押しながら使うと、マウス対応アプリでも
ローカルの文字選択を優先します。

\key{OSC 52} により、SSH 経由を含むターミナルのプログラムが Inkline の RAM
クリップボードを読み書きできます。ノートのクリップボードとは別です。
右 Alt/Opt+X はコピーを行い、削除はエディター内で行うよう案内します。
OSC 52 にはエディターの文字を削除する機能はありません。

\key{画像：} kitty のインライン／共有メモリー転送は RAM 内で処理します。
ファイル／一時ファイル転送は無効です。Sixel は明示的に要求すれば使えますが、
対応機能として自動通知はしません。kitty のプレースホルダー、アニメーション再生、
一部の sixel モードは未完成です。

\code{clear} と画面全体の消去は、文字と表示中の画像を消します。
メモリー不足時は画面外の画像を解放します。ディスクへの画像キャッシュはありませんが、
大きな画像は本体の限られた RAM を使います。

\topic{キーでプログラムを起動する}
Inkline のシェルで登録します。引数はそのまま渡されるので、シェル構文にはシェルを明示してください。
起動したプログラムの環境は Bash 経由で \code{~/.bashrc} を読み込みます。

\begin{shell}
~/inkline shortcut register e emacs
~/inkline shortcut register p python3
~/inkline shortcut list
~/inkline shortcut deregister p
\end{shell}

\key{Ctrl+Alt+E} で空きスロットに Emacs を開きます。
\code{~/inkline run PROGRAM ARG...} なら登録せずに起動できます。
再登録は既存の割り当てを置き換え、登録解除は動作中のプログラムには影響しません。
キーは US 配列の物理位置です。句読点は引用符で囲み、自作プログラムはフルパスを使います。
プログラムは事前にインストールしてください。

自分で画面を描画する Qt／電子ペーパー用アプリには次を使います。
\begin{shell}
~/inkline shortcut register a --epaper /home/root/my-app
\end{shell}
ネイティブアプリは同時に1つで、実行中は Inkline とシェルを一時停止します。
ランチャーは Qt の表示設定と RAM キャッシュを用意し、ログを破棄し、終了後に Inkline を再開します。
アプリはキーボードを排他的に取得せず、Qt の入力を使ってください。

\begin{notice}
\key{Ctrl+Alt+T は固定です。} 登録も解除もできません。
ネイティブアプリを終了して既存のターミナルに戻ります。
物理 Ctrl キーを使ってください。Caps Lock の置き換えは Inkline 内だけに適用されます。
\end{notice}

\key{緊急時：} Ctrl+Alt+Backspace を2秒間押すと、ショートカットアプリと子プロセスを停止し、
Inkline を再起動します。全ターミナルが閉じ、未保存の作業は失われます。
root アプリの入力独占、サービス変更、資源の枯渇により復旧できない場合は、本体を再起動します。

\topic{インストール・更新・マニュアル}
対象は \key{reMarkable 2、ファームウェア3.27} と標準 Qt 6.8 電子ペーパーライブラリーです。
最新の配布物とチェックサムの確認方法はこちらです。

\InstallLink

インストーラーは機種、ファームウェア、フォント、Qt 起動、RAM の一時領域、swap の無効化を確認します。
リリースは \path{/home/root/.local/share/inkline} 以下に置き、キーのランチャーは起動時に有効になります。
更新は開いているターミナルを維持します。\key{新しい版を使うには終了して開き直してください。}
旧リリースは切り戻し用に残るため、ストレージを使います。

この1つの PDF に9言語すべてを収録し、配布物に同梱します。
ノート画面が動作し、USB ウェブインターフェースが有効なら、インストール時に
\key{My files} への取り込みを試みます。成功した PDF は記録し、同じ版の重複取り込みを避けます。

自動取り込みができない場合は USB を接続し、Inkline を終了し、ノートのストレージ設定で
\key{USB web interface} を有効にしてから、SSH で実行します。
\begin{shell}
~/inkline manual
\end{shell}
PDF をダウンロードし、USB のブラウザー画面や reMarkable のデスクトップ／モバイルアプリから
取り込むこともできます。取り込み処理はノートのデータベースを直接変更せず、
ターミナルを中断しません。Inkline を削除しても、取り込んだ文書は My files に残ります。

SSH からノートへ戻る、または Inkline を削除するには次を使います。
\begin{shell}
~/inkline stop
~/inkline uninstall
\end{shell}
推奨 LaTeX 環境を含む追加ソフトウェアの導入手順はウェブサイトにあります。

\topic{仕組み・ストレージ・RAM}
Qt がキー入力を共通の変換処理に渡し、ショートカットと Caps Lock の動作を適用します。
libghostty-vt がターミナル用のキー列を生成し、各スロットの疑似端末（PTY）がプログラムと接続します。
シェルには \code{HOME=/home/root} を設定し、対話型 Bash は \code{~/.bashrc} を読みます。

出力はターミナルの解析器と sixel デコーダーを通ります。libghostty が文字セル、カーソル、
代替画面、履歴、kitty 画像を管理します。描画器はグレースケールの画像を保持して変更行だけを描き直し、
標準 Qt 電子ペーパーバックエンドが画面を更新します。独立した evdev デーモンが全体の
ショートカットと電源キーを処理し、systemd がスリープと復帰を担当します。

各ターミナルのコア割り当ては64 MiB、kitty 画像は画面ごとに16 MiB、保持する sixel 画像は
16 MiB が上限です。描画とアプリのメモリーは別です。空スロットはターミナルを確保しません。
大きなアプリ9個が同時に RAM に収まるとは限りません。

一時画像、キャッシュ、セッション状態は RAM を使います。インストール済みソフトウェア、
保存文書、終了時の設定、アプリが書くファイルは永続ストレージを使います。
Inkline は一般のプログラムによるファイル書き込みを禁止しません。

追加パッケージの Python は標準版です。キャッシュ書き込みを減らすには、本体の
\code{~/.bashrc} に次を追加します。
\begin{shell}
export PYTHONDONTWRITEBYTECODE=1
export PIP_NO_CACHE_DIR=1
\end{shell}
\code{source ~/.bashrc} を実行するか、新しいシェルを開きます。
pip のインストール時コンパイルも避けるには
\code{python3 -m pip install --no-compile PACKAGE} を使います。
Python の分離モードは Python 環境変数を無視します。詳細はウェブの Python ガイドにあります。

Inkline は GPL-3.0-or-later です。配布物には対応するソースと依存ソフトウェアの表記を含みます。
テスト記録では自動検証と、実機での確認が必要な動作を区別しています。
